% Add "answers" option to print the answers
\documentclass[addpoints]{exam}

\usepackage[utf8]{inputenc}
\usepackage{calc}
\usepackage[dvipsnames]{xcolor}
\usepackage{ifthen}
\usepackage{tcolorbox}

% \makebox[\textwidth]{Name:\enspace\hrulefill}
%
% \vspace{5mm}
%\makebox[\textwidth]{Instructor’s name:\enspace\hrulefill}

\definecolor{MyLightGray}{gray}{0.60}

\newcommand{\truefalse}
{		\begin{checkboxes}
			 \choice True
			 \choice False. Explanation:
		\end{checkboxes}
		\vspace{10pt}
}


% Check if 'answers' is an option
\newcommand{\ifanswers}[2]{%
  \ifthenelse{\boolean{exam@answers}}{#1}{#2}%
}



\newcommand{\answerspace}[1]
{
	\ifprintanswers
	\else
		\setlength{\mylength}{10pt}%
		% \vspace{#1\mylength}
		% Make the grey box
		\par\nobreak\noindent
		\makebox[\linewidth][l]{%
			% set inner padding to zero, so only the frame shows
			{\setlength{\fboxsep}{0pt}%
			% draw a color frame, white (or transparent) inside
			\fcolorbox{MyLightGray}{white}{%
			% the inner box has zero width content but fixed height
			\parbox[t][#1\mylength][c]{%
				% subtract twice the frame‐thickness from the full width
				\dimexpr\linewidth-2\fboxrule\relax
			}{\strut}%
			}%
			}%
		}%
		\par
		\vspace{3\mylength}

  	\fi
}
\newlength{\mylength}

\newcommand{\code}[1]{{\texttt{#1}}}


\usepackage{color}
\usepackage{lastpage}
\usepackage{verbatim}
\usepackage{fancyvrb}
\usepackage{graphicx}
\usepackage{float}
\usepackage{minted}
\usepackage{multicol}

% For those sweet, sweet acmart fonts.
\usepackage[T1]{fontenc}
\usepackage[tt=false, type1=true]{libertine}
\usepackage[varqu]{zi4}                  % Inconsolata for \ttfamily

\RecustomVerbatimCommand{\VerbatimInput}{VerbatimInput}{}

\newcommand{\Code}[1]{\mintinline{python}{#1}}

\renewcommand{\solutiontitle}{\noindent\textbf{Solution sketch: }}
\definecolor{SolutionColor}{rgb}{0.8,0.9,1}

\cfoot{Page \thepage\ of \pageref{LastPage}}
\rhead{Initials: \rule{20mm}{0.4pt}}

%%%%%%%%%%%%%%%%%%%%%%%%%%%%%%%%%%%%%%%%
\begin{document}
\fvset{numbers=left,numbersep=3pt}
\thispagestyle{empty}

\begin{center}
	{\LARGE CS 408 Fall 2026 -- Project 1 Quiz \\}
\fbox{\fbox{\parbox{5.5in}{\centering
      \begin{itemize}
        \item This quiz is closed notes, closed book.

        \item This quiz should be \textbf{completed individually}: any
          communication with your fellow students is \textbf{strictly}
          prohibited.

        \item If you have a question, raise your hand and the
          course staff will come to you.

        \item On multiple choice questions, draw a clear
          \textsf{circle} around your chosen answer.

        \item On open-ended questions, please print or write legibly.

        \end{itemize}
	}}}
\end{center}

\vspace{10mm}
\makebox[\textwidth]{\textbf{Name (please print)}: \rule{100mm}{0.8pt}\hfill}

 \vspace{10mm}
\makebox[\textwidth]{\textbf{PUID}: \rule{100mm}{0.8pt} \hfill}

\begin{questions}

  \vspace{12pt}
%% ------------------------------------------------------------
\question[1] Why do we need a harness to fuzz a library such as
  \Code{parson}?
  \begin{choices}
  \item The harness makes the library run faster.
    \CorrectChoice The harness takes the input from the fuzzer and feeds
    it to the fuzzing target.
  \item The harness decides which inputs to keep.
  \item The harness mutates the inputs for the fuzzer.
  \end{choices}
  \begin{solution}
    The harness is a bridge between the fuzzer and the target. The
    fuzzer creates the inputs and checks the results.
  \end{solution}

  \vspace{12pt}
%% ------------------------------------------------------------
\question[1] Suppose there are the following three branches in our
  fuzzing target \Code{project/vulnerable/vulnerable.c}:
  \begin{itemize}
  \item If the input does not start with \Code{"head "}, it prints
    \texttt{Input does not start with 'head '.} and returns.
  \item If the string length is greater than 8, it prints
    \texttt{Crash: string length greater than 8.} and crashes.
  \item If the number is greater than the string length, it prints
    \texttt{Hung: number is greater than string length.} and sleeps.
  \end{itemize}

  You start your fuzzer, which runs the target program on exactly
  these four inputs:

  \begin{center}
    \Code{"hello"} \quad\quad \Code{"head abcdefghij 1"} \quad\quad
    \Code{"head abcdefghijkl 2"} \quad\quad \Code{"head abc 9"}
  \end{center}

  What does the \Code{bugs=} field show in the line that your fuzzer
  prints every 5 seconds?
  \begin{choices}
  \item \Code{bugs=3(4)}
  \item \Code{bugs=1(3)}
  \CorrectChoice \Code{bugs=2(3)}
  \item \Code{bugs=3(3)}
  \end{choices}
  \begin{solution}
    For \Code{hello}, the target prints an error message but returns
    normally; the handout says this is not a bug. The two long strings
    crash with the same output, so they are \emph{one} unique bug.
    \Code{head abc 9} runs for more than 5 seconds, so it is a hang,
    and its output is different, so it is a second unique bug. Three
    bugs in total, two of them unique.
  \end{solution}

  \vspace{12pt}
%% ------------------------------------------------------------
\question[1] In Task 2, you ran your fuzzer with the \Code{--greybox}
  flag. Suppose your fuzzer runs the target program on a seed that
  covers no new edges. What happens?
  \begin{choices}
  \CorrectChoice The seed is not added to the corpus.
  \item The seed is added to the corpus, and its \Code{prev} is
    \Code{-1}.
  \item The seed is added to the corpus, and its \Code{prev} is
    \Code{0}.
  \item The fuzzer stops with an error.
  \end{choices}
  \begin{solution}
    The fuzzer runs each seed first, and adds it to the corpus only if
    it triggers new edge coverage. A seed in the corpus has
    \Code{prev} set to \Code{-1}.
  \end{solution}

\newpage

%% ------------------------------------------------------------
\question[1] Coverage collection has three parts: the compiler
  instrumentation, the instrumentation runtime (\Code{sancov.c}), and
  the fuzzer. What does each part do?
  \begin{choices}
  \item Instrumentation records the visited edges. The runtime decides
    which inputs to keep. The fuzzer adds the callbacks.
    \CorrectChoice Instrumentation inserts a callback when an edge is
    visited. The runtime updates the coverage map. The fuzzer reads
    the map and checks for new edges.
  \item The fuzzer adds the callbacks at runtime. The runtime mutates
    the inputs. Instrumentation reads the map.
  \item Instrumentation adds the callbacks. The runtime checks for new
    edges. The fuzzer only counts executions.
  \end{choices}
  \begin{solution}
    The harness sees only one execution. The fuzzer keeps track of all
    edges visited so far, so only the fuzzer can tell whether an edge
    is new.
  \end{solution}

  \vspace{12pt}
%% ------------------------------------------------------------
\question[1] In Task 2, you implemented
  \Code{__sanitizer_cov_trace_pc_guard} in \Code{sancov.c}. Suppose
  the guard of the visited edge holds the edge ID $k$. Which byte of
  the coverage map do you set to \Code{1}?
  \begin{choices}
  \item Byte $k$
  \item Byte $k - 1$
  \CorrectChoice Byte $4 + k - 1$
  \item Byte $4 + k$
  \end{choices}
  \begin{solution}
    The first four bytes store the total number of edges. Edge IDs
    start from 1, so edge $k$ uses byte $4 + k - 1$.
  \end{solution}

  \vspace{12pt}
%% ------------------------------------------------------------
\question[1] When your fuzzer discovers an input that triggers new
  code coverage, it minimizes that input by attempting to remove
  individual bytes from the input, one at a time. How does the fuzzer
  decide when to trim each byte?
  \begin{choices}
  \item When the input without that byte still covers every edge that
    the old input covered.
  \item When the input without that byte covers at least one new edge.
  \item When the input without that byte still crashes.
    \CorrectChoice When the input without that byte still triggers all
    the new code coverage.
  \end{choices}
  \begin{solution}
    This is the rule in the handout. The first option is too strict:
    almost no byte could ever be removed.
  \end{solution}

  \vspace{12pt}
%% ------------------------------------------------------------
\question[2] In Task 3, you ran your fuzzer with \Code{-g
    <grammar_file>}. How does your fuzzer mutate inputs from the
  corpus?
  \begin{choices}
  \item It applies the mutators from Task 1 to the input string.
  \item It generates a new input from the start nonterminal.
    \CorrectChoice It walks the rule decision tree for the selected
    input and expands some of the nonterminals using a different rule.
  \item It inserts tokens from the AFL dictionary into the string.
  \end{choices}
  \begin{solution}
    From the handout: ``Mutating an input is done by choosing a
    different production rule for a nonterminal.'' The handout also
    says not to use the Task 1 mutators in this task.
  \end{solution}

  \vspace{12pt}
%% ------------------------------------------------------------
\question[1] In Task 4, which part of the framework from the previous
  tasks does \Code{afl-cc} replace?
  \begin{choices}
  \item The harness and the seeds.
  \item The source code of the target library.
  \item The mutation and corpus logic in \Code{fuzzer.py}.
  \CorrectChoice The instrumentation flags and runtime.
  \end{choices}
  \begin{solution}
    The handout says you don't need to specify the instrumentation
    flags or implement the runtime yourself, because \Code{afl-cc}
    does it. The mutation and corpus logic is the job of
    \Code{afl-fuzz}.
  \end{solution}

\newpage

%% ------------------------------------------------------------
\question[3] Briefly describe how your team allocated responsibilities
  to team members. In particular, what parts of the project did you
  work on?  What tasks were your teamates primarily responsible for?
  How did your team coordinate tasks?  \answerspace{30} %
  \begin{solution}
    Any reasonable description of the student's own contribution and
    how the team divided the work.
  \end{solution}

\end{questions}

\end{document}
